
\FigureLayoutDeclare{img/2020/shanghai/q10_fig1.png}{6.0cm}{8bp 6bp 6bp 8bp}
\FigureLayoutDeclare{img/2020/shanghai/q12_fig1.png}{8.0cm}{9bp 6bp 6bp 7bp}

\chapter{2020年上海卷（秋）}

\section{填空题}

\begin{problem}
已知集合 \(A=\{1,2,4\}\)，集合 \(B=\{2,4,5\}\)，则 \(A\cap B=\)\fillinblank{}.
\end{problem}

\begin{answer}
\(\{2,4\}\)
\end{answer}

\begin{solution}
因为 \(A=\{1,2,4\}\)，\(B=\{2,4,5\}\)，两集合的公共元素为 \(2\)，\(4\)，所以 \(A\cap B=\{2,4\}\)，故答案为 \(\{2,4\}\).
\end{solution}

\begin{problem}
计算：\(\lim\limits_{n\to\infty}\frac{n+1}{3n-1}=\)\fillinblank{}.
\end{problem}

\begin{answer}
\(\frac13\)
\end{answer}

\begin{solution}
因为 \(n\ne 0\)，分子、分母同除以 \(n\)，得 \(\frac{n+1}{3n-1}=\frac{1+\frac1n}{3-\frac1n}\)，而 \(\frac1n\to 0\)，所以 \(\lim\limits_{n\to\infty}\frac{n+1}{3n-1}=\lim\limits_{n\to\infty}\frac{1+\frac1n}{3-\frac1n}=\frac13\)，故答案为 \(\frac13\).
\end{solution}

\begin{problem}
已知复数 \(z=1-2\i\)（\(\i\) 为虚数单位），则 \(|z|=\)\fillinblank{}.
\end{problem}

\begin{answer}
\(\sqrt5\)
\end{answer}

\begin{solution}
由 \(z=1-2\i\)，得 \(|z|=\sqrt{1^2+(-2)^2}=\sqrt5\)，故答案为 \(\sqrt5\).
\end{solution}

\begin{problem}
已知函数 \(f(x)=x^3\)，\(f^{-1}(x)\) 是 \(f(x)\) 的反函数，则 \(f^{-1}(x)=\)\fillinblank{}.
\end{problem}

\begin{answer}
\(\sqrt[3]{x}\)
\end{answer}

\begin{solution}
设 \(y=x^3\)，\(x\in\R\)，解得 \(x=\sqrt[3]{y}\)，把 \(x\) 与 \(y\) 互换，得反函数为 \(f^{-1}(x)=\sqrt[3]{x}\)，\(x\in\R\)，回代满足 \(f\left(f^{-1}(x)\right)=x\)，故答案为 \(\sqrt[3]{x}\).
\end{solution}

\begin{problem}
已知 \(x\)，\(y\) 满足 \(x+y-2\ge0\)，\(x+2y-3\le0\)，\(y\ge0\)，
则 \(z=y-2x\) 的最大值为\fillinblank{}.

\bitmapfigure[max width=6.0cm]{img/2020/shanghai/q5_fig1.png}
\end{problem}

\begin{answer}
\(-1\)
\end{answer}

\begin{solution}
约束 \(x+y-2\ge 0\)，\(x+2y-3\le 0\)，\(y\ge 0\) 的可行域是以 \((2,0)\)，\((3,0)\)，\(A(1,1)\) 为顶点的三角形及其内部，其中 \(A(1,1)\) 由 \(x+y-2=0\)，\(x+2y-3=0\) 联立解得，化目标函数 \(z=y-2x\) 为 \(y=2x+z\)，\(z\) 为直线在 \(y\) 轴上的截距，代入三个顶点得 \(z\) 的值分别为 \(-4\)，\(-6\)，\(-1\)，比较得最大值为 \(-1\)，在 \(A(1,1)\) 处取得，故答案为 \(-1\).
\end{solution}

\begin{problem}
已知行列式 \(\left|\begin{matrix}
1 & a & b\\
2 & c & d\\
3 & 0 & 0
\end{matrix}\right|=6\)，则 \(\left|\begin{matrix}
a & b\\
c & d
\end{matrix}\right|\)=\fillinblank{}.
\end{problem}

\begin{answer}
\(2\)
\end{answer}

\begin{solution}
按第三行展开，\(\left|\begin{matrix}1 & a & b\\ 2 & c & d\\ 3 & 0 & 0\end{matrix}\right|=1\cdot\left(c\cdot 0-d\cdot 0\right)-a\cdot\left(2\cdot 0-d\cdot 3\right)+b\cdot\left(2\cdot 0-c\cdot 3\right)=3ad-3bc=3\left(ad-bc\right)=6\)，所以 \(ad-bc=2\)，即 \(\left|\begin{matrix}a & b\\ c & d\end{matrix}\right|=2\)，故答案为 \(2\).
\end{solution}

\begin{problem}
已知有四个数 \(1\)，\(2\)，\(a\)，\(b\)，这四个数的中位数是 \(3\)，平均数是 \(4\)，则 \(ab=\)\fillinblank{}.
\end{problem}

\begin{answer}
\(36\)
\end{answer}

\begin{solution}
四个数的平均数为 \(4\)，所以和为 \(16\)，得 \(a+b=16-1-2=13\)，若 \(a\)，\(b\) 中有一个不大于 \(2\)，不妨设 \(a\le 2\)，则 \(b=13-a\ge 11\)，四个数按从小到大排列为 \(a\)，\(1\)，\(2\)，\(b\) 或 \(1\)，\(a\)，\(2\)，\(b\)，中位数为 \(\frac{1+2}{2}\) 或 \(\frac{a+2}{2}\)，前者为 \(\frac32\ne 3\)，后者等于 \(3\) 得 \(a=4\)，与 \(a\le 2\) 矛盾，故 \(a\)，\(b\) 均大于 \(2\)，此时排序为 \(1\)，\(2\) 及 \(a\)，\(b\)，中位数为 \(\frac{2+\min\left(a,b\right)}{2}=3\)，得 \(\min\left(a,b\right)=4\)，结合 \(a+b=13\) 得另一数为 \(9\)，回代排序后中间两数确为 \(2\)，\(4\)，中位数为 \(3\)，符合题意，所以 \(ab=4\times 9=36\)，故答案为 \(36\).
\end{solution}

\begin{problem}
已知数列 \(\{a_n\}\) 是公差不为零的等差数列，且 \(a_1+a_{10}=a_9\)，则
\(\frac{a_1+a_2+\cdots+a_9}{a_{10}}\)=\fillinblank{}.
\end{problem}

\begin{answer}
\(\frac{27}{8}\)
\end{answer}

\begin{solution}
设公差为 \(d\)，\(d\ne 0\)，则 \(a_n=a_1+\left(n-1\right)d\)，由 \(a_1+a_{10}=a_9\) 得 \(a_1+a_1+9d=a_1+8d\)，解得 \(a_1=-d\)，前 \(9\) 项和（即 \(S_9\)）为 \(a_1+a_2+\cdots+a_9=9a_1+36d=27d\)，\(a_{10}=a_1+9d=8d\ne 0\)，所以 \(\frac{a_1+a_2+\cdots+a_9}{a_{10}}=\frac{27d}{8d}=\frac{27}{8}\)，故答案为 \(\frac{27}{8}\).
\end{solution}

\begin{problem}
从 \(6\) 个人挑选 \(4\) 个人去值班，每人值班一天，第一天安排 \(1\) 个人，第二天安排 \(1\) 个人，第三天安排 \(2\) 个人，则共有\fillinblank{} 种安排情况.
\end{problem}

\begin{answer}
\(180\)
\end{answer}

\begin{solution}
先从\(6\)人中选\(1\)人排第一天，有\(\mathrm C_6^1\)种，再从余下\(5\)人中选\(1\)人排第二天，有\(\mathrm C_5^1\)种，最后从余下\(4\)人中选\(2\)人排第三天，有\(\mathrm C_4^2\)种，故不同安排方法共有\(\mathrm C_6^1\mathrm C_5^1\mathrm C_4^2=6\times5\times6=180\)种，故答案为\(180\).
\end{solution}

\begin{problem}
已知椭圆 \(C:\frac{x^2}{4}+\frac{y^2}{3}=1\) 的右焦点为 \(F\)，直线 \(l\) 经过椭圆右焦点 \(F\)，交椭圆 \(C\) 于 \(P\)，\(Q\) 两点（点 \(P\) 在第二象限），若点 \(Q\) 关于 \(x\) 轴对称点为 \(Q'\)，且满足 \(PQ\perp FQ'\)，求直线 \(l\) 的方程是\fillinblank{}.

\bitmapfigure[max width=6.0cm]{img/2020/shanghai/q10_fig1.png}
\end{problem}

\begin{answer}
\(x+y-1=0\)
\end{answer}

\begin{solution}
由\(a^2=4\)，\(b^2=3\)得\(c=\sqrt{4-3}=1\)，故右焦点为\(F(1,0)\).设直线\(PQ\)过\(F\)交椭圆于\(P\)，\(Q\)，记\(Q(x_Q,y_Q)\)，则\(Q^\prime(x_Q,-y_Q)\)亦在椭圆上.因为\(F\)在椭圆内，故\(\overrightarrow{FQ}\ne\bs{0}\)，且\(P\)，\(F\)，\(Q\)共线，所以\(FQ^\prime\perp FP\)等价于\(\overrightarrow{FQ}\cdot\overrightarrow{FQ^\prime}=0\)，即\((x_Q-1)^2-y_Q^2=0\)，得\(|x_Q-1|=|y_Q|\)，竖直直线\(x=1\)此时给出\(y_Q=0\)，不可能，故直线\(PQ\)必有斜率.联立\(\frac{x_Q^2}{4}+\frac{y_Q^2}{3}=1\)得\(7x_Q^2-8x_Q-8=0\)，故过\(F\)的直线斜率为\(1\)或\(-1\).直线\(y=x-1\)与椭圆的两个交点为\(\left(\frac{4+6\sqrt2}{7},\frac{-3+6\sqrt2}{7}\right)\)，\(\left(\frac{4-6\sqrt2}{7},\frac{-3-6\sqrt2}{7}\right)\)，分别位于第四、第三象限，不含第二象限点，舍去.直线\(y=1-x\)与椭圆的两个交点为\(\left(\frac{4+6\sqrt2}{7},\frac{3-6\sqrt2}{7}\right)\)，\(\left(\frac{4-6\sqrt2}{7},\frac{3+6\sqrt2}{7}\right)\)，后者横坐标为负、纵坐标为正，位于第二象限，可取为\(P\)，符合题意.故直线\(PQ\)的方程为\(x+y-1=0\).（原 PDF 记椭圆为 \(\Gamma\)，垂直条件为 \(FQ'\perp FP\)，在 \(P\)，\(F\)，\(Q\) 共线时与现稿等价.）
\end{solution}

\begin{problem}
设 \(a\in\R\)，若存在定义域为 \(\R\) 的函数 \(f(x)\) 同时满足下列两个条件：

\begin{enumerate}
\item 对任意的 \(x_0\in\R\)，\(f(x_0)\) 的值为 \(x_0\) 或 \(x_0^2\)；
\item 关于 \(x\) 的方程 \(f(x)=a\) 无实数解，
\end{enumerate}

则 \(a\) 的取值范围是\fillinblank{}.
\end{problem}

\begin{answer}
\((-\infty,0)\cup(0,1)\cup(1,+\infty)\)
\end{answer}

\begin{solution}
对 \(x_0=0\)，\(\{x_0,x_0^2\}=\{0\}\)，故必有 \(f(0)=0\)，对 \(x_0=1\)，\(\{x_0,x_0^2\}=\{1\}\)，故必有 \(f(1)=1\)，所以 \(0\) 与 \(1\) 恒在 \(f\) 的值域中，方程 \(f(x)=a\) 无实数解蕴含 \(a\ne0\) 且 \(a\ne1\).反之，对任意 \(a\in\R\) 且 \(a\ne0\)，\(a\ne1\)，定义 \(f(a)=a^2\)，其余 \(x\) 处取 \(f(x)=x\)，则每点取值确为 \(x\) 或 \(x^2\)；\(f(x)=x\) 取到 \(a\) 只可能在 \(x=a\) 处，但该处已改为 \(a^2\ne a\)，\(f(a)=a^2=a\) 将推出 \(a=0\) 或 \(a=1\)，已排除，故 \(a\) 不在值域中.因此满足条件的 \(a\) 恰为 \((-\infty,0)\cup(0,1)\cup(1,+\infty)\).
\end{solution}

\begin{problem}
已知 \(\bs{a}_1\)，\(\bs{a}_2\)，\(\bs{b}_1\)，\(\bs{b}_2\)，\(\ldots\)，\(\bs{b}_k\)（\(k\in\N^*\)）是平面内两两互不相等的向量，满足 \(|\bs{a}_1-\bs{a}_2|=1\)，且 \(|\bs{a}_i-\bs{b}_j|\in\{1,2\}\)（其中 \(i=1\)，\(2\)，\(j=1\)，\(2\)，\(\ldots\)，\(k\)），则 \(k\) 的最大值是\fillinblank{}.

\bitmapfigure[max width=6.0cm]{img/2020/shanghai/q12_fig1.png}
\end{problem}

\begin{answer}
\(6\)
\end{answer}

\begin{solution}
设\(\overrightarrow{OA_1}=\bs{a}_1\)，\(\overrightarrow{OA_2}=\bs{a}_2\)，则\(|A_1A_2|=1\).每个\(\bs{b}_j\)对应的点\(B_j\)到\(A_1\)，\(A_2\)的距离都属于\(\{1,2\}\)，故必为圆\(C(A_1,r_1)\)与\(C(A_2,r_2)\)（\(r_1,r_2\in\{1,2\}\)）的交点.圆心距\(d=1\)，\(r_1=r_2=1\)时\(0<1<2\)，相交于\(2\)点，\(r_1=r_2=2\)时\(0<1<4\)，相交于\(2\)点，\(\{r_1,r_2\}=\{1,2\}\)时\(|r_1-r_2|=1=d\)，内切于\(1\)点.四组交点中，到\(A_1\)距离不同的两组不可能相交，到\(A_1\)距离同为\(1\)的切点落在直线\(A_1A_2\)上，而\((1,1)\)的两交点构成等边三角形而不落在该直线上，故六点两两不同，且它们到\(A_1\)，\(A_2\)的距离非零，故与\(\bs{a}_1\)，\(\bs{a}_2\)及彼此都互不相等.于是\(k\)至多为\(2+1+1+2=6\)，且取这六点时取到，故\(k\)的最大值是\(6\).
\end{solution}
\section{单选题}

\begin{problem}
下列不等式恒成立的是（\quad）

\choices
    {\(a^2+b^2\le2ab\)}
    {\(a^2+b^2\ge-2ab\)}
    {\(a+b\ge2\sqrt{|ab|}\)}
    {\(a^2+b^2\le-2ab\)}
\end{problem}

\begin{answer}
B
\end{answer}

\begin{solution}
A.取 \(a=1\)，\(b=-1\)，则 \(a^2+b^2=2\)，\(2ab=-2\)，不等式 \(a^2+b^2\le2ab\) 不成立，故 A 错误；B.因为 \((a+b)^2\ge0\)，即 \(a^2+b^2+2ab\ge0\)，所以 \(a^2+b^2\ge-2ab\) 恒成立，故 B 正确；C.取 \(a=b=-1\)，则 \(a+b=-2\)，\(2\sqrt{|ab|}=2\)，不等式不成立，故 C 错误；D.按本页现稿取 \(a=b=1\)，则 \(a^2+b^2=2\)，\(-2ab=-2\)，不等式 \(a^2+b^2\le-2ab\) 不成立，故 D 错误.原 PDF 的选项 D 为 \(a+b\ge-2\sqrt{|ab|}\)，取 \(a=-4\)，\(b=-1\) 得 \(-5\ge-4\) 亦不成立，故按两种题面结论同为 B.故选 B.
\end{solution}

\begin{problem}
已知直线方程 \(3x+4y+1=0\) 的一个参数方程可以是（\quad）

\choices
    {\(x=1+3t\)，\(y=-1-4t\)（\(t\) 为参数）}
    {\(x=1-4t\)，\(y=-1+3t\)（\(t\) 为参数）}
    {\(x=1-3t\)，\(y=-1+4t\)（\(t\) 为参数）}
    {\(x=1+4t\)，\(y=1-3t\)（\(t\) 为参数）}
\end{problem}

\begin{answer}
B
\end{answer}

\begin{solution}
按本页现稿：点 \((1,-1)\) 满足 \(3\times1+4\times(-1)+1=0\)，在直线 \(3x+4y+1=0\) 上.方向向量 \((-4,3)\) 满足 \(3\times(-4)+4\times3=0\)，与直线平行，故 \(x=1-4t\)，\(y=-1+3t\)（\(t\) 为参数）是该直线的一个参数方程.A 的方向向量为 \((3,-4)\)，\(3\times3+4\times(-4)\ne0\)，C 的方向向量为 \((-3,4)\)，\(3\times(-3)+4\times4\ne0\)，D 的点 \((1,1)\) 不在直线上，均不合要求.故选 B.原 PDF 四个选项为 A：\(x=1+3t\)，\(y=-1+4t\)；B：\(x=1-3t\)，\(y=-1+4t\)；C：\(x=1-4t\)，\(y=-1-3t\)；D：\(x=1+4t\)，\(y=-1-3t\)（\(t\) 为参数），其中只有 D 的方向向量 \((4,-3)\) 与直线平行且过点 \((1,-1)\)，故按原 PDF 题面应选 D.
\end{solution}

\begin{problem}
在棱长为 \(10\) 的正方体 \(ABCD-A_1B_1C_1D_1\) 中，\(P\) 为左侧面 \(ADD_1A_1\) 上一点，已知点 \(P\) 到 \(A_1D_1\) 的距离为 \(3\)，\(P\) 到 \(AA_1\) 的距离为 \(2\)，则过点 \(P\) 且与 \(A_1C\) 平行的直线交正方体于 \(P\)，\(Q\) 两点，则 \(Q\) 点所在的平面是（\quad）

\bitmapfigure[max width=6.0cm]{img/2020/shanghai/q15_fig1.png}

\choices
    {\(AA_1B_1B\)}
    {\(BB_1C_1C\)}
    {\(CC_1D_1D\)}
    {\(ABCD\)}
\end{problem}

\begin{answer}
D
\end{answer}

\begin{solution}
建立空间直角坐标系，设\(A(0,0,0)\)，\(B(10,0,0)\)，\(C(10,10,0)\)，\(D(0,10,0)\)，\(A_1(0,0,10)\)，\(C_1(10,10,10)\).点\(P\)在侧面\(ADD_1A_1\)即平面\(x=0\)上，记\(P(0,y,z)\).\(P\)到直线\(AA_1\)（即\(x=0\)，\(y=0\)）的距离为\(|y|=2\)，得\(y=2\)（\(P\) 在面内，舍去 \(y=-2\)），到直线\(A_1D_1\)（即\(x=0\)，\(z=10\)）的距离为\(|10-z|=3\)，得\(z=7\)（舍去 \(z=13\)），故\(P(0,2,7)\).\(\overrightarrow{A_1C}=(10,10,-10)\)平行于\((1,1,-1)\)，过\(P\)且平行于\(A_1C\)的直线为\((t,2+t,7-t)\)，\(t\in\R\).当\(t=7\)时得点\((7,9,0)\)，\(x\)，\(y\in[0,10]\) 且 \(z=0\)，落在底面\(ABCD\)上；\(t<0\) 时 \(x<0\) 已在体外，\(t>7\) 时 \(z<0\) 已在体外，故另一交点\(Q\)即在此处.所以点\(Q\)所在的面是\(ABCD\)，按本页现稿为选项 D，故选 D；原 PDF 同一内容为选项 A.
\end{solution}

\begin{problem}
命题 \(p\)：存在 \(a\in\R\) 且 \(a\neq0\)，对于任意的 \(x\in\R\)，使得 \(f(x+a)<f(x)+f(a)\)；

命题 \(q_1\)：\(f(x)\) 单调递减且 \(f(x)>0\) 恒成立；

命题 \(q_2\)：\(f(x)\) 单调递增，存在 \(x_0<0\) 使得 \(f(x_0)=0\)，

则下列说法正确的是（\quad）

\choices
    {只有 \(q_1\) 是 \(p\) 的充分条件}
    {只有 \(q_2\) 是 \(p\) 的充分条件}
    {\(q_1\)，\(q_2\) 都是 \(p\) 的充分条件}
    {\(q_1\)，\(q_2\) 都不是 \(p\) 的充分条件}
\end{problem}

\begin{answer}
C
\end{answer}

\begin{solution}
若\(q_1\)成立，任取\(a>0\)（则\(a\ne0\)），对任意\(x\in\R\)，由\(x+a>x\)及\(f\)单调递减得\(f(x+a)<f(x)\)，又\(f(a)>0\)，故\(f(x)<f(x)+f(a)\)，于是\(f(x+a)<f(x)+f(a)\)对一切\(x\)成立，即\(q_1\)是\(p\)的充分条件.若\(q_2\)成立，取\(a=x_0<0\)（则\(a\ne0\)），对任意\(x\in\R\)，由\(x+a<x\)及\(f\)单调递增得\(f(x+a)<f(x)\)，又\(f(a)=f(x_0)=0\)，故\(f(x)+f(a)=f(x)>f(x+a)\)对一切\(x\)成立，即\(q_2\)也是\(p\)的充分条件.因此\(q_1\)，\(q_2\)都是\(p\)的充分条件，按本页现稿为选项 C，故选 C；原 PDF 同一内容为选项 A.
\end{solution}
\section{解答题}

\begin{problem}
已知 \(ABCD\) 是边长为 \(1\) 的正方形，正方形 \(ABCD\) 绕 \(AB\) 旋转形成一个圆柱.

\begin{enumerate}
\item 求该圆柱的表面积；
\item 正方形 \(ABCD\) 绕 \(AB\) 逆时针旋转 \(\frac{\pi}{2}\) 至 \(ABC_1D_1\)，求线段 \(CD_1\) 与平面 \(ABCD\) 所成的角.
\end{enumerate}

\bitmapfigure[max width=6.0cm]{img/2020/shanghai/q17_fig1.png}
\end{problem}

\begin{solution}
\begin{enumerate}
\item 正方形边长为 \(1\)，绕 \(AB\) 旋转时以 \(AB\) 为轴，点 \(C\)，\(D\) 到轴的距离均为 \(1\)，故圆柱底面半径为 \(1\)，高为 \(AB=1\)，表面积 \(S=2\pi\times1^2+2\pi\times1\times1=4\pi\).
\item 记原正方形所在平面为 \(\alpha\)，旋转轴 \(AB\) 上的点不动，\(D\) 在与轴垂直的平面内转过 \(\frac{\pi}{2}\) 到 \(D_1\)，则 \(AD_1=AD=1\)，\(AD_1\perp AD\)，且 \(AD_1\) 仍与轴 \(AB\) 垂直，由 \(AB\cap AD=A\) 得 \(AD_1\perp\alpha\)，即 \(D_1\) 在 \(\alpha\) 上的投影为 \(A\)，于是 \(AC\) 为 \(CD_1\) 在 \(\alpha\) 上的投影，所求角为 \(\angle D_1CA\)，由 \(AC=\sqrt2\)，\(AD_1=1\) 得 \(CD_1=\sqrt{AC^2+AD_1^2}=\sqrt3\)，\(\cos\angle D_1CA=\frac{AC}{CD_1}=\frac{\sqrt6}{3}\)，故所成角的大小为 \(\arccos\frac{\sqrt6}{3}\).
\end{enumerate}（原 PDF 作绕 \(BC\) 旋转至 \(A_1BCD_1\)、求 \(AD_1\) 与平面 \(ABCD\) 所成角，对称地结论同为 \(4\pi\) 与 \(\arccos\frac{\sqrt6}{3}\).）
\end{solution}

\begin{problem}
已知函数 \(f(x)=\sin\omega x\)，\(\omega>0\).

\begin{enumerate}
\item \(f(x)\) 的周期是 \(4\pi\)，求 \(\omega\)，并求 \(f(x)=\frac12\) 的解集；
\item 已知 \(\omega=1\)，\(g(x)=f^2(x)+\sqrt3\,f(-x)f\!\left(\frac\pi2-x\right)\)，\(x\in\left[0,\frac\pi4\right]\)，求 \(g(x)\) 的值域.
\end{enumerate}
\end{problem}

\begin{solution}
\begin{enumerate}
\item 由周期 \(T=\frac{2\pi}{\omega}=4\pi\) 得 \(\omega=\frac12\)，即 \(f(x)=\sin\frac{x}{2}\)，由 \(\sin\frac{x}{2}=\frac12\) 得 \(\frac{x}{2}=2k\pi+\frac{\pi}{6}\) 或 \(\frac{x}{2}=2k\pi+\frac{5\pi}{6}\)，\(k\in\Z\)，故解集为 \(\{4k\pi+\frac{\pi}{3},4k\pi+\frac{5\pi}{3}\mid k\in\Z\}\).
\item 当 \(\omega=1\) 时 \(f(x)=\sin x\)，故 \(g(x)=\sin^2x+\sqrt3\sin(-x)\sin\left(\frac{\pi}{2}-x\right)=\frac{1-\cos2x}{2}-\frac{\sqrt3}{2}\sin2x=\frac12-\sin\left(2x+\frac{\pi}{6}\right)\)，由 \(x\in\left[0,\frac{\pi}{4}\right]\) 得 \(2x+\frac{\pi}{6}\in\left[\frac{\pi}{6},\frac{2\pi}{3}\right]\)，\(\sin\left(2x+\frac{\pi}{6}\right)\in\left[\frac12,1\right]\)，故 \(g(x)\in\left[-\frac12,0\right]\)，即值域为 \(\left[-\frac12,0\right]\).
\end{enumerate}
\end{solution}

\begin{problem}
在研究某市交通情况时，道路密度是指该路段上一定时间内通过的车辆数除以时间，车辆密度是该路段一定时间内通过的车辆数除以该路段的长度，现定义交通流量 \(v=f(x)=\frac{q}{x}\)，其中 \(x\) 为道路密度，\(q\) 为车辆密度，且
\(f(x)=100-135\left(\frac13\right)^{\frac{80}{x}}\)（\(0<x<40\)），
\(f(x)=-k(x-40)+85\)（\(40\le x\le80\)）.

\begin{enumerate}
\item 若交通流量 \(v>95\)，求道路密度 \(x\) 的取值范围；
\item 已知道路密度 \(x=80\) 时，测得交通流量 \(v=50\)，求车辆密度 \(q\) 的最大值.
\end{enumerate}
\end{problem}

\begin{solution}
\begin{enumerate}
\item 原 PDF 注明 \(k>0\)（本页现稿遗漏），以下按 \(k>0\) 求解；实测 \(k=\frac78>0\)，一致.当 \(40\le x\le80\) 时 \(f(x)=-k(x-40)+85\le85<95\)，故 \(v>95\) 只可能在 \(0<x<40\) 内，由 \(100-135\left(\frac13\right)^{\frac{80}{x}}>95\) 得 \(\left(\frac13\right)^{\frac{80}{x}}<\frac1{27}\)，即 \(\frac{80}{x}>3\)，亦即 \(0<x<\frac{80}{3}\)，且 \(\frac{80}{3}<40\) 自动成立，故所求取值范围为 \(\left(0,\frac{80}{3}\right)\).
\item 把 \(x=80\)，\(v=50\) 代入 \(v=-k(x-40)+85\) 得 \(50=-40k+85\)，解得 \(k=\frac78\)，于是
\[q=vx=
\begin{cases}
x\left(100-135\left(\frac13\right)^{\frac{80}{x}}\right), & 0<x<40,\\
\left(-\frac78(x-40)+85\right)x, & 40\le x\le80.
\end{cases}\]
当 \(0<x<40\) 时 \(v<100\)，故 \(q<4000\)，当 \(40\le x\le80\) 时 \(q=-\frac78x^2+120x\)，对称轴 \(x=\frac{480}{7}\in[40,80]\)，最大值为 \(\frac{28800}{7}>4000\)，故 \(q\) 的最大值为 \(\frac{28800}{7}\).
\end{enumerate}
\end{solution}

\begin{problem}
已知双曲线 \(\Gamma_1:\frac{x^2}{4}-\frac{y^2}{b^2}=1\) 与圆 \(\Gamma_2:x^2+y^2=4+b^2\)（\(b>0\)）交于点 \(A(x_A,y_A)\)（第一象限），曲线 \(\Gamma\) 为 \(\Gamma_1\)，\(\Gamma_2\) 上取满足 \(x>|x_A|\) 的部分.

\begin{enumerate}
\item 若 \(x_A=\sqrt6\)，求 \(b\) 的值；
\item 当 \(b=\sqrt5\)，\(\Gamma_2\) 与 \(x\) 轴交点记作点 \(F_1\)，\(F_2\)，\(P\) 是曲线 \(\Gamma\) 上一点，且在第一象限，且 \(|PF_1|=8\)，求 \(\angle F_1PF_2\)；
\item 过点 \(D\left(0,\frac{b^2}{2}+2\right)\) 斜率为 \(-\frac{b}{2}\) 的直线 \(l\) 与曲线 \(\Gamma\) 只有两个交点，记为 \(M\)，\(N\)，用 \(b\) 表示 \(\overrightarrow{OM}\cdot\overrightarrow{ON}\)，并求 \(\overrightarrow{OM}\cdot\overrightarrow{ON}\) 的取值范围.
\end{enumerate}
\end{problem}

\begin{solution}
\begin{enumerate}
\item 点 \(A\) 在第一象限且同时满足 \(\frac{x_A^2}{4}-\frac{y_A^2}{b^2}=1\)，\(x_A^2+y_A^2=4+b^2\)，代入 \(x_A=\sqrt6\) 得 \(\frac{y_A^2}{b^2}=\frac12\)，\(y_A^2=b^2-2\)，由 \(y_A>0\)，\(b>0\) 解得 \(y_A=\sqrt2\)，\(b=2\).
\item 当 \(b=\sqrt5\) 时 \(\Gamma\) 与 \(x\) 轴的交点为 \(F_1(-3,0)\)，\(F_2(3,0)\)，恰为双曲线的两焦点（本页现稿称 \(F_1\)，\(F_2\) 为 \(\Gamma_2\) 的焦点，原 PDF 称为 \(\Gamma\) 与 \(x\) 轴交点，此处两者一致），由联立方程得 \(y_A^2=\frac{25}{9}\)，\(x_A^2=\frac{56}{9}\)，圆半径为 \(3\)，圆上任一点 \(P\) 满足 \(|PF_1|\le|PO|+|OF_1|=6<8\)，故满足 \(|PF_1|=8\) 的第一象限点 \(P\) 必在双曲线右支 \(x\ge x_A\) 部分上，由双曲线定义 \(|PF_1|-|PF_2|=4\) 得 \(|PF_2|=4\)，又 \(|F_1F_2|=6\)，由余弦定理 \(\cos\angle F_1PF_2=\frac{64+16-36}{2\times8\times4}=\frac{11}{16}\)，故 \(\angle F_1PF_2=\arccos\frac{11}{16}\).
\item 本页现稿过点记作 \(D\)（原 PDF 记作 \(S\)，坐标同为 \(\left(0,\frac{b^2}{2}+2\right)\)），记 \(t=\frac{b^2+4}{2}\)，直线 \(l\) 为 \(y=-\frac{b}{2}x+t\)，原点到 \(l\) 的距离为 \(\frac{t}{\sqrt{1+b^2/4}}=\sqrt{4+b^2}\)，恰为圆半径，故 \(l\) 与圆相切，切点满足 \(y=\frac{2}{b}x\)，解得切点 \(M(b,2)\)，\(l\) 斜率与渐近线 \(y=-\frac{b}{2}x\) 平行，代入双曲线后二次项相消得关于 \(x\) 的一次方程 \(\frac{t}{b}x-\frac{t^2}{b^2}-1=0\)，有唯一交点 \(N\)，横坐标 \(x_N=\frac{t}{b}+\frac{b}{t}>0\)，由 \(OM\perp l\) 及 \(N\) 在 \(l\) 上得 \(\overrightarrow{OM}\cdot\overrightarrow{ON}=|\overrightarrow{OM}|^2=4+b^2\)，联立 \(\Gamma_1\)，\(\Gamma_2\) 得 \(y_A^2=\frac{b^4}{4+b^2}\)，\(M\) 与 \(A\) 同在圆 \(x^2+y^2=4+b^2\) 上，故 \(b\ge x_A\) 等价于 \(y_A\ge2\)，等价于 \(b^4\ge4(4+b^2)\)，即 \(b^2\ge2+2\sqrt5\)，另记 \(s=b^2\)，则 \(x_N^2-x_A^2=\frac{(s^2+12s+16)^2-32s(s+4)(s+2)}{4s(s+4)^2}=\frac{(s^2-4s-16)^2}{4s(s+4)^2}\ge0\)，故 \(x_N\ge x_A\) 恒成立，等号仅当 \(s=2+2\sqrt5\)（另一根 \(2-2\sqrt5<0\) 因 \(s>0\) 不取），此时 \(M=N=A\)，于是 \(l\) 与 \(\Gamma\) 恰有两个不同公共点当且仅当 \(b^2>2+2\sqrt5\)，此时 \(\overrightarrow{OM}\cdot\overrightarrow{ON}=4+b^2\)，取值范围为 \((6+2\sqrt5,+\infty)\).
\end{enumerate}（原 PDF 记双曲线为 \(C_1\)、圆为 \(C_2\)，曲线 \(\Gamma\) 取 \(|x|\ge x_A\) 部分；等号情形 \(M=N=A\) 为单点，两种取法下恰有两个不同公共点当且仅当 \(b^2>2+2\sqrt5\)，结论一致.）
\end{solution}

\begin{problem}
已知数列 \(\{a_n\}\) 为有限数列，满足 \(|a_1-a_2|\le|a_1-a_3|\le\cdots\le|a_1-a_m|\)，则称 \(\{a_n\}\) 满足性质 \(P\).

\begin{enumerate}
\item 判断数列 \(3\)，\(2\)，\(5\)，\(1\) 和 \(4\)，\(3\)，\(2\)，\(5\)，\(1\) 是否具有性质 \(P\)，请说明理由；
\item 若 \(a_1=1\)，公比为 \(q\) 的等比数列，项数为 \(10\)，具有性质 \(P\)，求 \(q\) 的取值范围；
\item 若 \(\{a_n\}\) 是 \(1\)，\(2\)，\(3\)，\(\ldots\)，\(m\) 的一个排列（\(m\ge4\)），\(\{b_n\}\) 符合 \(b_k=a_{k+1}\)（\(k=1\)，\(2\)，\(\ldots\)，\(m-1\)），\(\{a_n\}\)，\(\{b_n\}\) 都具有性质 \(P\)，求所有满足条件的数列 \(\{a_n\}\).
\end{enumerate}
\end{problem}

\begin{solution}
\begin{enumerate}
\item 数列 \(3\)，\(2\)，\(5\)，\(1\) 的各距离为 \(|2-3|=1\)，\(|5-3|=2\)，\(|1-3|=2\)，满足 \(1\le2\le2\)，具有性质 \(P\)，数列 \(4\)，\(3\)，\(2\)，\(5\)，\(1\) 的各距离为 \(1\)，\(2\)，\(1\)，\(3\)，因 \(|5-4|=1<2=|2-4|\)，不具有性质 \(P\).
\item 记 \(a_n=q^{n-1}\)，\(n=1\)，\(\ldots\)，\(10\)，性质 \(P\) 即 \(|q^j-1|\ge|q^{j-1}-1|\) 对 \(j=1\)，\(\ldots\)，\(9\) 成立，平方相减得 \((q-1)q^{j-1}[q^{j-1}(q+1)-2]\ge0\)，\(j=1\) 时恒成立，当 \(q=0\) 时数列为 \(1\)，\(0\)，\(\ldots\)，\(0\)，距离恒为 \(1\)，满足性质 \(P\)，当 \(q\ge1\) 时 \(q^{j-1}(q+1)\) 随 \(j\) 递增，只需 \(j=1\) 处 \(q+1-2\ge0\)，恒成立，当 \(0<q\le1\) 时 \(q^{j-1}(q+1)\) 随 \(j\) 递减，只需 \(j=1\) 处 \(q+1-2\le0\)，恒成立，故一切 \(q\ge0\) 满足性质 \(P\)，当 \(-1\le q<0\) 时 \(j=2\) 处 \(|q^2-1|=1-q^2\)，\(|q-1|=1-q\)，\(1-q^2\ge1-q\) 等价于 \(q(q-1)\le0\)，而此时 \(q(q-1)>0\)，不成立，当 \(q<-1\) 时 \(j\) 为奇数则 \(q^{j-1}>0\)，\(q^{j-1}(q+1)-2<0\) 恒成立，\(j\) 为偶数则需 \(q^{j-1}(q+1)-2\ge0\)，左端等于 \(|q|^{j-1}|q+1|\) 随偶数 \(j\) 递增，只需 \(j=2\) 处 \(q(q+1)-2=(q+2)(q-1)\ge0\)，即 \(q\le-2\)，充分性由单调性保证，故 \(q\) 的取值范围为 \((-\infty,-2]\cup[0,+\infty)\).
\item 设 \(a_1=p\)（本页现稿 \(b_k=a_{k+1}\)；原 PDF 作 \(b_k=a_{k-1}\)，\(k=1\) 时下标越界，以下按现稿求解），若 \(p=1\)，则距离集为 \(\{1,\ldots,m-1\}\) 各恰一次，非减即迫使 \(a_i=i\)，得 \(1\)，\(2\)，\(\ldots\)，\(m\)，其 \(b\) 为 \(2\)，\(\ldots\)，\(m\)，距离 \(1\)，\(\ldots\)，\(m-2\) 非减，满足条件，若 \(p=m\)，同理距离集为 \(\{1,\ldots,m-1\}\)，迫使 \(a\) 为 \(m\)，\(m-1\)，\(\ldots\)，\(1\)，其 \(b\) 距离 \(1\)，\(\ldots\)，\(m-2\) 非减，满足条件，若 \(p=2\)，则距离多重集为 \(\{1,1,2,\ldots,m-2\}\)，\(a_2\)，\(a_3\) 必为 \(1\)，\(3\) 的某种排列，其后必为 \(4\)，\(\ldots\)，\(m\)，若 \(a_2=3\)，则 \(b_1=3\)，\(|a_3-b_1|=2\)，\(|a_4-b_1|=1<2\)，\(b\) 不具性质 \(P\)，故只剩 \(2\)，\(1\)，\(3\)，\(\ldots\)，\(m\)，其 \(b\) 距离为 \(2\)，\(3\)，\(\ldots\)，\(m-1\)，满足条件，若 \(p=m-1\)，对称地 \(a_2\)，\(a_3\) 必为 \(m\)，\(m-2\) 的某种排列，其后必为 \(m-3\)，\(\ldots\)，\(1\)，若 \(a_2=m-2\)，则 \(|a_3-b_1|=2\)，\(|a_4-b_1|=1<2\)，\(b\) 不具性质 \(P\)，故只剩 \(m-1\)，\(m\)，\(m-2\)，\(\ldots\)，\(1\)，满足条件，若 \(3\le p\le m-2\)（\(m=4\) 时此情形为空），则 \(p\pm1\)，\(p\pm2\) 均在 \(1\)，\(\ldots\)，\(m\) 内，\(a_2\)，\(a_3\) 必为 \(p-1\)，\(p+1\) 的某种排列，\(a_4\)，\(a_5\) 必为 \(p-2\)，\(p+2\) 的某种排列，四种组合下 \(b_1=a_2\) 的距离依次为：\(2\)，\(1\)；\(2\)，\(3\)，\(1\)；\(2\)，\(3\)，\(1\)；\(2\)，\(1\)，其中每组都有一次后项严格小于前项（依次为 \(1<2\)；\(1<3\)；\(1<3\)；\(1<2\)），故 \(b\) 均不具性质 \(P\)，于是满足条件的数列恰为以上四种：\(1\)，\(2\)，\(\ldots\)，\(m\)；\(m\)，\(m-1\)，\(\ldots\)，\(1\)；\(2\)，\(1\)，\(3\)，\(\ldots\)，\(m\)；\(m-1\)，\(m\)，\(m-2\)，\(\ldots\)，\(1\).
\end{enumerate}
\end{solution}
